\begin{abstract}
Koopman modeling represents nonlinear dynamics as linear evolution in a lifted space, but its accuracy depends critically on the lifting map and on how the linear operator is identified. Most existing approaches construct observables from instantaneous states and identify a single time-invariant operator, so they neglect short-term temporal correlations and cannot distinguish how the system evolves from different past trajectories. We propose Transformer-Koopman, a temporal lifting framework in which a patch-based Transformer encoder maps a finite history window to a structured lifted state, and the lifted transition is history-conditioned by a low-rank modulation of a full-rank nominal operator. We show that a one-step-trained history-conditioned operator develops a training--deployment inconsistency under recursive scheduled propagation that spectral stabilization does not repair, and that training under the same scheduled rollout used at deployment restores long-horizon accuracy. The predictor is integrated into a one-iteration real-time-iteration (RTI) MPC that avoids re-running the full encoder at every prediction stage, achieving a $p95$ controller latency of 120\,ms, within the sampling period, while tracking a matched exact scheduled NMPC oracle within $1\%$. A finite-horizon theory characterizes the rollout-error propagation and the local RTI approximation, and empirical diagnostics assess its assumptions. Experiments on a 6-state manipulator and a 2-state soft robotic platform support the modeling and real-time control claims.
\end{abstract}
